% SOURCE_AUTHORITY: sga5_fr_workpass.tex, lines 8033--8097 (LF-normalized unit).
% SOURCE_SHA256: 791F4EFFC5E02832D5D77ED03518C8156D6F07E4C8238B03545DB93D883FBB28
Suponhamos doravante que \(X\) é estrito. As seguintes expressões para \(Z_r^{pq}(X_n)\) e \(B_r^{pq}(X_n)\):
\[
Z_r^{pq}(X_n)=
\Im\!\left[
T^{p+q-1}\bigl(F^{p-r+1}X_n/F^pX_n\bigr)
\xrightarrow{\partial}
T^{p+q}\bigl(F^pX_n/F^{p+1}X_n\bigr)
\right],
\]
\[
B_r^{pq}(X_n)=
\Im\!\left[
T^{p+q}\bigl(F^pX_n/F^{p+r}X_n\bigr)
\longrightarrow
T^{p+q}\bigl(F^pX_n/F^{p+1}X_n\bigr)
\right],
\]
mostram que
\begin{equation}
\left\{
\begin{array}{ll}
Z_r^{pq}(X_m)\xrightarrow{\sim}Z_r^{pq}(X_n)&\text{para }m\ge n\ge p,\\[3pt]
B_r^{pq}(X_m)\xrightarrow{\sim}B_r^{pq}(X_n)&\text{para }m\ge n\ge p+r-1,
\end{array}
\right.
\tag{\(d_1\)}
\end{equation}
onde os morfismos canônicos são isomorfismos. Em particular, o morfismo canônico
\[
\mathcal E_r^{pq}(X)\longrightarrow E_r^{pq}(X_{p+r-1})
\]
é um isomorfismo. Além disso, de b) segue-se que
\begin{equation}
\left\{
\begin{array}{ll}
Z_\infty^{pq}(X_{p+r-1})=Z_r^{pq}(X_{p+r-1})&(r\ge1),\\[3pt]
B_\infty^{pq}(X_{p+r-1})=B_r^{pq}(X_{p+r-1})&(r>p).
\end{array}
\right.
\tag{\(d_2\)}
\end{equation}
Isto permite definir, para todo \(r\), um epimorfismo canônico
\[
\theta_r^{pq}:\mathcal E_r^{pq}(X)\longrightarrow E_\infty^{pq}(X_{p+r-1})
\]
que é um isomorfismo quando \(r>p\). Além disso, se \(a\) é um inteiro tal que
\[
\mathcal B_r^{pq}(X)=\mathcal B_a^{pq}(X)\qquad(r\ge a),
\]
Por exemplo, podemos tomar \(a=p+1\), para que
\begin{equation}
\cdots\hookrightarrow\mathcal E_r^{pq}(X)\hookrightarrow\cdots\hookrightarrow
\mathcal E_{a+1}^{pq}(X)\hookrightarrow\mathcal E_a^{pq}(X),
\tag{S bis}
\end{equation}
Então, para \(s\ge r\ge a\), os diagramas
\begin{equation}
\begin{array}{ccc}
\mathcal E_s^{pq}(X)&\xrightarrow[\epi]{\theta_s^{pq}}&E_\infty^{pq}(X_{p+s-1})\\
\downarrow&&\downarrow\can\\
\mathcal E_r^{pq}(X)&\xrightarrow[\epi]{\theta_r^{pq}}&E_\infty^{pq}(X_{p+r-1})
\end{array}
\tag{\(d_3\)}
\end{equation}
são comutativos.
